\documentclass[12pt,a4paper]{article}
\usepackage{../latex/dsdlab-report}

\begin{document}
	\reporttitlepage{۱۰}{پیاده‌سازی یک پردازنده‌ی ساده}

	\section*{۱۰.۱ مقدمه و اهداف}
	هدف این آزمایش، پیاده‌سازی یک پردازنده با معماری پشته‌ای\LTRfootnote{Stack Architecture} و سپس توسعهٔ آن است. مطابق دستورکار، این پردازنده پشته‌ای با هشت ثبات هشت‌بیتی دارد، مجموعه دستورش هشت دستور است، و حافظه‌ای به اندازهٔ ۲۵۶ خانهٔ هشت‌بیتی دارد که هشت خانهٔ آخر آن، یعنی آدرس‌های \lr{F8} تا \lr{FF}، به‌صورت ورودی/خروجی نگاشته‌شده در حافظه\LTRfootnote{Memory-Mapped I/O} به کار می‌روند.

	افزون بر مجموعه دستور پایه، این آزمایش سه دستور تازه نیز می‌خواهد و برنامه‌ای که یک تابع چندبخشی\LTRfootnote{Piecewise Function} را محاسبه کند. تابع خواسته‌شده چنین است:
	\[
	Y =
	\begin{cases}
	3X + 5,  & 0 \le X < 32,\\
	2X + 20, & 32 \le X < 64,\\
	X - 40,  & 64 \le X \le 127.
	\end{cases}
	\]
	مقدار \(X\) از ورودی برد خوانده و مقدار \(Y\) روی نمایشگرهای هفت‌قسمتی\LTRfootnote{Seven Segment Display} نشان داده می‌شود. اگر ورودی منفی باشد یا خروجی از محدودهٔ قابل نمایش بیرون بزند، یکی از چراغ‌های برد به نشانهٔ خطا روشن می‌گردد.

	\subsection*{روند انجام آزمایش}
	ابتدا مجموعه دستور پایه و سه دستور تازه مشخّص و ریزعملیات\LTRfootnote{Micro-operation} هر یک استخراج شد. سپس پردازنده به‌صورت یک ماشین حالت متناهی\LTRfootnote{Finite State Machine} نوشته شد و در کنار آن، حافظه و درگاه‌های ورودی/خروجی در یک پیمانهٔ سامانه قرار گرفتند. برنامهٔ زبان ماشین با یک همگردان\LTRfootnote{Assembler} کوچک به کد ماشین تبدیل شد و درستی آن پیش از هر شبیه‌سازی سخت‌افزاری، با یک شبیه‌ساز مستقل برای هر ۲۵۶ ورودی ممکن وارسی گردید. سرانجام سه میز آزمون\LTRfootnote{Testbench} در \lr{ModelSim} اجرا و طرح در \lr{Quartus II} سنتز و تحلیل زمانی شد.

	\section*{۱۰.۲ مجموعه دستور}
	\subsection*{الف) دستورهای پایه}
	هشت دستور پایه همان‌هایی هستند که دستورکار تعریف کرده است. کد عملیات\LTRfootnote{Opcode} در چهار بیت پرارزش بایت دستور قرار می‌گیرد.

	\begin{table}[H]
		\centering
		\footnotesize
		\begin{tabular}{|c|l|c|l|}
			\hline
			\textbf{کد} & \textbf{دستور} & \textbf{طول} & \textbf{عملکرد} \\ \hline
			\(0000\) & \lr{PUSHC C} & \(2\) & مقدار ثابت \lr{C} را در پشته می‌گذارد \\ \hline
			\(0001\) & \lr{PUSH M}  & \(2\) & محتوای خانه یا درگاه \lr{M} را در پشته می‌گذارد \\ \hline
			\(0010\) & \lr{POP M}   & \(2\) & از پشته برمی‌دارد و در خانه یا درگاه \lr{M} می‌نویسد \\ \hline
			\(0011\) & \lr{JUMP}    & \(1\) & از پشته برمی‌دارد و در شمارندهٔ برنامه می‌گذارد \\ \hline
			\(0100\) & \lr{JZ}      & \(1\) & اگر پرچم \lr{Z} یک باشد، همان کار \lr{JUMP} \\ \hline
			\(0101\) & \lr{JS}      & \(1\) & اگر پرچم \lr{S} یک باشد، همان کار \lr{JUMP} \\ \hline
			\(0110\) & \lr{ADD}     & \(1\) & دو مقدار بالای پشته را جمع و نتیجه را جایگزین می‌کند \\ \hline
			\(0111\) & \lr{SUB}     & \(1\) & دو مقدار بالای پشته را تفریق و نتیجه را جایگزین می‌کند \\ \hline
		\end{tabular}
		\caption{هشت دستور پایه‌ی خواسته‌شده در دستورکار}
	\end{table}

	دو قرارداد مهم دستورکار نیز رعایت شده‌اند: تنها دو پرچم \lr{S} و \lr{Z} وجود دارد و این دو تنها با اجرای \lr{ADD} و \lr{SUB} تغییر می‌کنند، و همهٔ محاسبات علامت‌دار و در نمایش متمّم دو\LTRfootnote{Two's Complement} انجام می‌شوند.

	\subsection*{ب) سه دستور تازه}
	سه دستور زیر به مجموعه افزوده شدند. هر سه تک‌بایتی‌اند و عملوند جداگانه‌ای ندارند.

	\begin{table}[H]
		\centering
		\footnotesize
		\begin{tabular}{|c|l|l|l|}
			\hline
			\textbf{کد} & \textbf{دستور} & \textbf{عملکرد} & \textbf{اثر بر پشته و پرچم‌ها} \\ \hline
			\(1000\) & \lr{SWAP} & جای دو مقدار بالای پشته را عوض می‌کند & عمق و پرچم‌ها بی‌تغییر \\ \hline
			\(1001\) & \lr{SHL}  & مقدار بالای پشته را یک بیت به چپ می‌راند & عمق و پرچم‌ها بی‌تغییر \\ \hline
			\(1010\) & \lr{CMP}  & \(B-A\) را می‌سنجد بدون آنکه چیزی بردارد & تنها پرچم‌ها به‌روز می‌شوند \\ \hline
		\end{tabular}
		\caption{سه دستور افزوده‌شده به مجموعه‌ی پایه}
	\end{table}

	انگیزهٔ هر سه دستور از نیازهای همان برنامهٔ خواسته‌شده می‌آید. \lr{SHL} برای ضرب در دو لازم است؛ چون مجموعهٔ پایه هیچ دستور ضربی ندارد و تابع خواسته‌شده در دو شاخهٔ خود به \(2X\) و \(3X\) نیاز دارد. با شیفت یک‌بیتی، ضرب در دو بدون هیچ بلوک ضرب‌کننده‌ای به دست می‌آید و \(3X\) نیز از جمع \(X\) با \(2X\) ساخته می‌شود.

	\lr{CMP} برای تشخیص بازه لازم است. با مجموعهٔ پایه، تنها راه مقایسه استفاده از \lr{SUB} است، امّا \lr{SUB} هر دو عملوند را از پشته برمی‌دارد و نتیجه را جای آن‌ها می‌گذارد؛ یعنی مقدار \(X\) که هنوز به آن نیاز داریم از میان می‌رود. \lr{CMP} همان تفریق را انجام می‌دهد و تنها پرچم‌ها را به‌روز می‌کند، بی‌آنکه محتوای پشته را دست بزند.

	\lr{SWAP} نیز برای مرتّب کردن عملوندها در پشته پیش‌بینی شده است. در برنامهٔ نهایی به آن نیازی نیفتاد، امّا چون بخشی از توسعهٔ خواسته‌شده است، پیاده‌سازی و جداگانه آزموده شد.

	\section*{۱۰.۳ نمودار بلوکی پردازنده‌ی توسعه‌یافته}
	پردازنده از شمارندهٔ برنامه، ثبات دستور، ثبات عملوند، پشتهٔ هشت‌تایی به‌همراه اشاره‌گر پشته، منطق محاسباتی، دو پرچم و واحد کنترل ساخته شده است. مسیر داده برای سه دستور تازه به همان پشته و همان مسیر جمع و تفریق موجود متّصل شده تا توسعه با کمترین سخت‌افزار اضافی انجام شود.

	\shot{screenshots/block_diagram.png}{نمودار بلوکی پردازنده‌ی توسعه‌یافته}

	اگر \(A\) را مقدار بالای پشته و \(B\) را مقدار دوم از بالا بگیریم، روابط اصلی مسیر داده چنین‌اند:
	\[
	A = \mathrm{Stack}[\mathrm{SP}-1], \qquad B = \mathrm{Stack}[\mathrm{SP}-2],
	\]
	\[
	\text{جمع} = B + A, \qquad
	\text{تفریق} = B - A, \qquad
	\text{شیفت} = \{A[6{:}0],\,0\}.
	\]
	توجّه شود که هم \lr{SUB} و هم \lr{CMP} مقدار \(B-A\) را می‌سازند، یعنی مقدار دوم منهای مقدار بالا. این ترتیب در نوشتن برنامه اهمّیّت دارد و در بخش ۱۰.۶ به آن بازمی‌گردیم.

	\section*{۱۰.۴ تغییرات مسیر داده برای سه دستور تازه}
	\shot{screenshots/datapath_new_instr.png}{اثر سه دستور تازه بر پشته و پرچم‌ها}

	\subsection*{الف) \lr{SWAP}}
	این دستور جای دو مقدار بالای پشته را عوض می‌کند. نکتهٔ پیاده‌سازی آن این است که هر دو انتساب در یک لبهٔ ساعت و به‌صورت غیرمسدودکننده\LTRfootnote{Non-Blocking Assignment} انجام می‌شوند؛ بنابراین سمت راست هر دو انتساب مقادیر \emph{پیش از} لبه را می‌خواند و هیچ‌یک از دو مقدار پیش از استفاده بازنویسی نمی‌شود. اگر این دو انتساب مسدودکننده بودند، مقدار نخست پیش از خوانده شدن از بین می‌رفت و به‌جای جابه‌جایی، هر دو خانه یک مقدار می‌گرفتند. نه اشاره‌گر پشته و نه پرچم‌ها تغییر نمی‌کنند.

	\subsection*{ب) \lr{SHL}}
	مقدار بالای پشته یک بیت به چپ رانده می‌شود؛ بیت کم‌ارزش تازه صفر است و بیت پرارزش پیشین دور ریخته می‌شود. عمق پشته و پرچم‌ها دست‌نخورده می‌مانند. دور ریخته شدن بیت پرارزش پیامد مهمی دارد که در تحلیل برنامه به آن می‌پردازیم: اگر \(2X\) از محدودهٔ هشت‌بیتی علامت‌دار بیرون بزند، بیت علامت نتیجه یک می‌شود و همین، سازوکار تشخیص خطا را فراهم می‌کند.

	\subsection*{ج) \lr{CMP}}
	مقدار موقّت \(T=B-A\) روی همان مسیر تفریق موجود ساخته می‌شود، پرچم \lr{Z} از صفر بودن \(T\) و پرچم \lr{S} از بیت هفتم آن به دست می‌آید، و خودِ \(T\) هیچ‌جا ذخیره نمی‌شود. اشاره‌گر پشته و محتوای آن نیز تغییر نمی‌کنند. این دستور تنها دستور مجموعه است که پرچم‌ها را بدون تغییر دادن پشته به‌روز می‌کند.

	\section*{۱۰.۵ واحد کنترل و ریزعملیات}
	واحد کنترل یک ماشین حالت متناهی با چهارده حالت است. چرخهٔ عمومی از واکشی\LTRfootnote{Fetch} به رمزگشایی\LTRfootnote{Decode} و سپس به حالت اجرایی مناسب می‌رود و از آنجا به واکشی بازمی‌گردد. برای سه دستور تازه، سه حالت اجرایی جداگانه افزوده شده است. چون هر سه تک‌بایتی‌اند، برخلاف \lr{PUSHC}، \lr{PUSH} و \lr{POP} وارد حالت واکشی عملوند نمی‌شوند.

	\shot{screenshots/fsm_diagram.png}{نمودار حالت واحد کنترل به‌همراه سه حالت تازه}

	\begin{table}[H]
		\centering
		\footnotesize
		\begin{tabular}{|l|l|l|}
			\hline
			\textbf{دستور} & \textbf{مرحله} & \textbf{ریزعملیات} \\ \hline
			هر سه & واکشی    & \(\mathrm{IR} \leftarrow M[\mathrm{PC}]\)؛ \(\mathrm{PC} \leftarrow \mathrm{PC}+1\) \\ \hline
			هر سه & رمزگشایی & چهار بیت پرارزش ثبات دستور، حالت اجرایی را تعیین می‌کند \\ \hline
			\lr{SWAP} & اجرا & \(\mathrm{Stack}[\mathrm{SP}-1] \leftarrow B\)؛ \(\mathrm{Stack}[\mathrm{SP}-2] \leftarrow A\) \\ \hline
			\lr{SHL}  & اجرا & \(\mathrm{Stack}[\mathrm{SP}-1] \leftarrow \{A[6{:}0],0\}\) \\ \hline
			\lr{CMP}  & اجرا & \(T \leftarrow B-A\)؛ \(Z \leftarrow (T=0)\)؛ \(S \leftarrow T[7]\) \\ \hline
			هر سه & پایان & اعلام پایان دستور و بازگشت به واکشی \\ \hline
		\end{tabular}
		\caption{ریزعملیات سه دستور تازه در هر مرحله}
	\end{table}

	واحد کنترل حالت‌های مرزی پشته را نیز مدیریت می‌کند. \lr{SWAP} و \lr{CMP} تنها هنگامی کار اصلی خود را انجام می‌دهند که دست‌کم دو مقدار در پشته باشد و \lr{SHL} تنها هنگامی که دست‌کم یک مقدار باشد. اگر عمق پشته کافی نباشد، ماشین بدون هیچ دسترسی نامعتبری به واکشی بازمی‌گردد. این وارسی از خواندن خانه‌های نامعیّن پشته جلوگیری می‌کند.

	\vcode{code/stack_cpu.v}{پردازنده‌ی پشته‌ای با سه دستور افزوده}

	\section*{۱۰.۶ برنامه‌ی زبان ماشین}
	\subsection*{الف) نگاشت حافظه و درگاه‌ها}
	ارتباط پردازنده با ورودی و خروجی، همان‌گونه که دستورکار تصریح کرده، تنها از راه ورودی/خروجی نگاشته‌شده در حافظه انجام می‌شود. یعنی برنامه هیچ دستور ویژهٔ ورودی/خروجی ندارد و با همان \lr{PUSH} و \lr{POP} معمولی به درگاه‌ها دسترسی پیدا می‌کند.

	\begin{table}[H]
		\centering
		\small
		\begin{tabular}{|c|l|l|}
			\hline
			\textbf{آدرس} & \textbf{نقش} & \textbf{جهت} \\ \hline
			\lr{F8} & مقدار ورودی \(X\) از کلیدهای برد & خواندنی \\ \hline
			\lr{F9} & رقم کم‌ارزش نمایشگر هفت‌قسمتی & نوشتنی \\ \hline
			\lr{FA} & رقم پرارزش نمایشگر هفت‌قسمتی & نوشتنی \\ \hline
			\lr{FB} & چراغ خطا & نوشتنی \\ \hline
			\lr{FC} & چراغ پایان کار & نوشتنی \\ \hline
		\end{tabular}
		\caption{نگاشت درگاه‌های ورودی/خروجی در فضای حافظه}
	\end{table}

	پنج خانهٔ \lr{E0} تا \lr{E4} نیز به‌عنوان متغیّرهای موقّت برنامه به کار می‌روند؛ به‌ترتیب برای نگه‌داشتن \(X\)، دور ریختن عملوندهای \lr{CMP}، نگه‌داشتن \(Y\)، باقی‌مانده، و شمارندهٔ رقم پرارزش. این خانه‌ها بیرون از محدودهٔ کد برنامه‌اند و با آن تداخل ندارند.

	\subsection*{ب) ساختار برنامه}
	\shotw{screenshots/flowchart.png}{0.5}{فلوچارت برنامه‌ی زبان ماشین}

	برنامه نخست \(X\) را از درگاه \lr{F8} می‌خواند و در \lr{E0} ذخیره می‌کند. سپس با سه مقایسهٔ پیاپی، بازهٔ ورودی تعیین می‌شود. ترتیب عملوندها در همهٔ مقایسه‌ها یکسان است: نخست \(X\) و سپس حدّ مرزی در پشته گذاشته می‌شوند، پس \(B=X\) و \(A=\)حد است و \lr{CMP} مقدار \(X-\)حد را می‌سنجد. با این ترتیب، پرچم \lr{S} دقیقاً هنگامی یک می‌شود که \(X\) کوچک‌تر از آن حد باشد و دستور \lr{JS} می‌تواند مستقیماً برای شاخه‌بندی به کار رود.

	\begin{table}[H]
		\centering
		\footnotesize
		\begin{tabular}{|c|l|l|}
			\hline
			\textbf{مقایسه} & \textbf{شرط} & \textbf{مقصد} \\ \hline
			\(X - 0\)  & \(S=1\) یعنی \(X<0\)  & مسیر خطا \\ \hline
			\(X - 32\) & \(S=1\) یعنی \(X<32\) & شاخهٔ نخست \\ \hline
			\(X - 64\) & \(S=1\) یعنی \(X<64\) & شاخهٔ دوم \\ \hline
			--- & در غیر این صورت & شاخهٔ سوم \\ \hline
		\end{tabular}
		\caption{سه مقایسه‌ی پیاپی و شاخه‌بندی برنامه}
	\end{table}

	نکتهٔ ظریف در همین جدول نهفته است. چون \lr{CMP} پشته را دست نمی‌زند، پس از هر مقایسه باید دو عملوند با دو \lr{POP} به خانهٔ دورریز \lr{E1} تخلیه شوند. اگر این تخلیه انجام نمی‌شد، پشتهٔ هشت‌تایی پس از چند مقایسه پر می‌شد. همین رفتار، دلیل وجود خانهٔ \lr{E1} در جدول متغیّرهاست.

	\subsection*{ج) محاسبه‌ی سه شاخه}
	هر شاخه با ترکیب دستورهای موجود ساخته می‌شود:
	\begin{itemize}
		\item \textbf{شاخهٔ نخست} \(3X+5\): دو نسخه از \(X\) روی پشته گذاشته می‌شود، یکی با \lr{SHL} به \(2X\) تبدیل می‌گردد و \lr{ADD} مقدار \(X+2X=3X\) را می‌سازد؛ سپس عدد ۵ افزوده می‌شود.
		\item \textbf{شاخهٔ دوم} \(2X+20\): \(X\) با \lr{SHL} به \(2X\) تبدیل و عدد ۲۰ به آن افزوده می‌شود.
		\item \textbf{شاخهٔ سوم} \(X-40\): عدد ۴۰ از \(X\) کم می‌شود.
	\end{itemize}

	\subsection*{د) تشخیص خطا}
	دو وضعیّت باید به روشن شدن چراغ خطا بینجامد. نخست، ورودی منفی: مقادیر \(128\) تا \(255\) در نمایش هشت‌بیتی علامت‌دار منفی‌اند و همان نخستین مقایسه با صفر آن‌ها را می‌گیرد. دوم، خروجی خارج از محدوده: این حالت تنها در شاخهٔ دوم رخ می‌دهد، چون برای \(X\) میان \(54\) و \(63\)، مقدار ریاضی \(2X+20\) میان \(128\) و \(146\) می‌افتد که از \(127\) بزرگ‌تر است. در چنین حالتی بیت علامت نتیجهٔ هشت‌بیتی یک می‌شود و یک \lr{JS} پس از محاسبه، اجرا را به مسیر خطا می‌برد.

	دو شاخهٔ دیگر هرگز سرریز نمی‌کنند و این را می‌توان با محاسبهٔ ساده نشان داد: بیشینهٔ شاخهٔ نخست در \(X=31\) برابر \(3(31)+5=98\) و بیشینهٔ شاخهٔ سوم در \(X=127\) برابر \(127-40=87\) است؛ هر دو کمتر از \(127\).

	\subsection*{ه) تولید دو رقم نمایشگر}
	مقدار \(Y\) باید به‌صورت دو رقم شانزده‌شانزدهی روی نمایشگرها برود. رقم کم‌ارزش آسان است: سخت‌افزار درگاه تنها چهار بیت کم‌ارزش دادهٔ نوشته‌شده را به نمایشگر می‌دهد، پس نوشتن خودِ \(Y\) در \lr{F9} کافی است. رقم پرارزش دشوارتر است، چون مجموعهٔ دستور نه دستور \lr{AND} دارد و نه شیفت به راست. راه‌حل به‌کاررفته، تفریق پیاپی است: در یک حلقه، مقدار ۱۶ از باقی‌مانده کم و شمارنده یکی زیاد می‌شود تا وقتی که باقی‌مانده منفی شود؛ شمارندهٔ نهایی همان \(\lfloor Y/16 \rfloor\) است.

	\vcode[language={}]{code/lab10_program.asm}{برنامه‌ی زبان ماشین به زبان اسمبلی}

	کد ماشین این برنامه با یک همگردان کوچک تولید می‌شود که در بستهٔ تحویلی آمده است. برنامه از آدرس \lr{00} آغاز و در \lr{9B} پایان می‌یابد. فهرست کامل کد ماشین، شامل آدرس، بایت‌ها و دستور متناظر، در همان بسته و در پروندهٔ فهرست موجود است.

	\section*{۱۰.۷ حداکثر عمق پشته}
	یکی از خواسته‌های گزارش، تعیین بیشینهٔ عمق پشته در طول اجرای برنامه است. اهمّیّت این عدد در آن است که پشتهٔ سخت‌افزاری تنها هشت خانه دارد و سرریز آن رفتار برنامه را خراب می‌کند.

	با بررسی برنامه می‌توان دید که در هیچ لحظه‌ای بیش از دو مقدار هم‌زمان روی پشته نیست: الگوی کار همواره «گذاشتن یک یا دو مقدار، انجام عملیات، برداشتن نتیجه» است و پیش از هر شاخه، پشته کاملاً تخلیه می‌شود. برای آنکه این ادّعا صرفاً استدلالی نماند، بیشینهٔ عمق در شبیه‌سازی نیز اندازه‌گیری شد و برای هر ۲۵۶ ورودی ممکن برابر \textbf{دو} به دست آمد؛ یعنی از هشت خانهٔ موجود تنها دو خانه به کار می‌رود و حاشیهٔ اطمینان چهاربرابری وجود دارد.

	\section*{۱۰.۸ وارسی مستقل برنامه پیش از شبیه‌سازی سخت‌افزار}
	پیش از آنکه برنامه روی توصیف سخت‌افزار اجرا شود، درستی خودِ برنامه با یک شبیه‌ساز مستقل وارسی شد. این شبیه‌ساز همان مجموعه دستور را در سطح رفتاری اجرا می‌کند و برنامه را برای \emph{هر ۲۵۶ ورودی ممکن} می‌آزماید.

	\docode[language={}]{code/verify_program.py}{شبیه‌ساز مستقل وارسی برنامه}

	ارزش این گام در جدا کردن دو گونه خطاست. اگر نتیجه‌ای نادرست باشد، این وارسی می‌گوید که مشکل از خودِ برنامه است یا از توصیف سخت‌افزار؛ چون اگر شبیه‌ساز مستقل درست جواب بدهد و سخت‌افزار ندهد، ایراد در توصیف سخت‌افزار است و برعکس. خروجی این وارسی، درستی هر ۲۵۶ ورودی و بیشینهٔ عمق پشتهٔ برابر دو را تأیید می‌کند.

	\section*{۱۰.۹ میزهای آزمون و اجرا در \lr{ModelSim}}
	سه میز آزمون خودآزما نوشته شد که هر یک لایهٔ متفاوتی از طرح را می‌سنجد.

	\begin{table}[H]
		\centering
		\footnotesize
		\begin{tabular}{|l|l|}
			\hline
			\textbf{میز آزمون} & \textbf{آنچه سنجیده می‌شود} \\ \hline
			دستورهای تازه & رفتار \lr{SWAP}، \lr{SHL} و \lr{CMP} به‌صورت جداگانه \\ \hline
			کنترل \lr{JZ}   & پرش شرطی بر پایه‌ی پرچم \lr{Z} \\ \hline
			برنامه‌ی نهایی & اجرای کامل برنامه برای ده ورودی نماینده \\ \hline
		\end{tabular}
		\caption{سه میز آزمون آزمایش}
	\end{table}

	\docode{code/run_all.do}{پرونده‌ی فرمان اجرای همه‌ی آزمون‌ها}

	\subsection*{الف) آزمون سه دستور تازه}
	\shotw{screenshots/sim_1.png}{0.78}{اجرای آزمون سه دستور تازه}
	\shotw{screenshots/sim_2.png}{0.75}{خروجی متنی آزمون سه دستور تازه}

	این آزمون هشت وارسی دارد. یک وارسی برای \lr{SWAP}، چهار وارسی برای \lr{SHL} با مقادیر مرزی، و سه وارسی برای \lr{CMP} در هر سه حالت ممکن مقایسه. جدول زیر نتیجهٔ هر یک را با دلیل آن می‌آورد.

	\begin{table}[H]
		\centering
		\footnotesize
		\begin{tabular}{|l|l|l|l|}
			\hline
			\textbf{دستور} & \textbf{ورودی} & \textbf{نتیجه} & \textbf{چرا} \\ \hline
			\lr{SHL} & \(0\)   & \(0\)   & شیفت صفر همچنان صفر است \\ \hline
			\lr{SHL} & \(1\)   & \(2\)   & ضرب در دو \\ \hline
			\lr{SHL} & \(63\)  & \(126\) & بزرگ‌ترین ورودی بدون سرریز \\ \hline
			\lr{SHL} & \(127\) & \(254\) & بیت پرارزش دور ریخته می‌شود؛ نتیجه در متمّم دو منفی است \\ \hline
			\lr{CMP} & \(B=5, A=7\) & \(S=1, Z=0\) & \(5-7<0\) \\ \hline
			\lr{CMP} & \(B=5, A=5\) & \(S=0, Z=1\) & \(5-5=0\) \\ \hline
			\lr{CMP} & \(B=7, A=5\) & \(S=0, Z=0\) & \(7-5>0\) \\ \hline
		\end{tabular}
		\caption{وارسی‌های آزمون سه دستور تازه و دلیل هر نتیجه}
	\end{table}

	وارسی چهارم \lr{SHL} از همه گویاتر است. ورودی \(127\) بزرگ‌ترین عدد مثبت هشت‌بیتی است و شیفت آن مقدار \(254\) می‌دهد که در متمّم دو برابر \(-2\) است. یعنی همان سازوکاری که در بخش ۱۰.۶ برای تشخیص سرریز شاخهٔ دوم به کار رفت، اینجا به‌صورت جداگانه آزموده شده است.

	\subsection*{ب) آزمون پرش شرطی}
	\shotw{screenshots/sim_3.png}{0.82}{شکل موج آزمون پرش شرطی \lr{JZ}}
	\shotw{screenshots/sim_4.png}{0.72}{خروجی متنی آزمون \lr{JZ}}

	این آزمون کوچک، درستی مسیر پرش شرطی را جدا از بقیهٔ طرح می‌سنجد؛ یعنی بررسی می‌کند که با یک بودن پرچم \lr{Z}، مقدار بالای پشته واقعاً در شمارندهٔ برنامه بنشیند و عمق پشته یک واحد کم شود.

	\subsection*{ج) آزمون برنامه‌ی نهایی}
	\shotw{screenshots/sim_5.png}{0.82}{شکل موج اجرای کامل برنامه}
	\shotw{screenshots/sim_6.png}{0.75}{خروجی متنی اجرای برنامه برای ده ورودی}

	ده ورودی این آزمون تصادفی نیستند و هر یک برای پوشاندن یک وضعیّت مشخّص انتخاب شده‌اند. جدول زیر همهٔ آن‌ها را با محاسبهٔ دستی و نتیجهٔ مشاهده‌شده می‌آورد.

	\begin{table}[H]
		\centering
		\footnotesize
		\begin{tabular}{|c|l|l|c|c|l|}
			\hline
			\(X\) & \textbf{شاخه} & \textbf{محاسبه} & \(Y\) & \textbf{نمایش} & \textbf{هدف وارسی} \\ \hline
			\(0\)   & نخست & \(3(0)+5\)    & \(5\)   & \(05\) & کمترین ورودی معتبر \\ \hline
			\(31\)  & نخست & \(3(31)+5\)   & \(98\)  & \(62\) & آخرین مقدار شاخهٔ نخست \\ \hline
			\(32\)  & دوم  & \(2(32)+20\)  & \(84\)  & \(54\) & نخستین مقدار شاخهٔ دوم \\ \hline
			\(53\)  & دوم  & \(2(53)+20\)  & \(126\) & \(7E\) & آخرین مقدار بدون سرریز \\ \hline
			\(54\)  & دوم  & \(2(54)+20\)  & \(128\) & خطا & نخستین سرریز \\ \hline
			\(63\)  & دوم  & \(2(63)+20\)  & \(146\) & خطا & آخرین مقدار شاخهٔ دوم \\ \hline
			\(64\)  & سوم  & \(64-40\)     & \(24\)  & \(18\) & نخستین مقدار شاخهٔ سوم \\ \hline
			\(127\) & سوم  & \(127-40\)    & \(87\)  & \(57\) & بیشترین ورودی معتبر \\ \hline
			\(128\) & ---  & ورودی منفی    & ---     & خطا & نخستین مقدار منفی \\ \hline
			\(255\) & ---  & ورودی منفی    & ---     & خطا & \(-1\) در متمّم دو \\ \hline
		\end{tabular}
		\caption{ده ورودی آزمون برنامه، محاسبه‌ی دستی و هدف هر وارسی}
	\end{table}

	انتخاب این ده مقدار پوشش کاملی می‌سازد. سه جفت \((31,32)\)، \((53,54)\) و \((63,64)\) دقیقاً روی مرزهای تصمیم قرار دارند و هر جفت، دو سوی یک مرز را می‌آزماید؛ اگر شرط مقایسه‌ای یک واحد جابه‌جا نوشته شده بود، دست‌کم یکی از این جفت‌ها آن را آشکار می‌کرد. دو مقدار \(128\) و \(255\) نیز مسیر ورودی منفی را می‌پوشانند.

	\shotw{screenshots/sim_7.png}{0.75}{ادامه‌ی خروجی اجرای برنامه}
	\shotw{screenshots/sim_8.png}{0.82}{شکل موج اجرای همه‌ی آزمون‌ها پشت سر هم}
	\shotw{screenshots/sim_9.png}{0.75}{گزارش پایانی اجرای همه‌ی آزمون‌ها}

	\subsection*{د) تحلیل شکل موج اجرای برنامه}
	در شکل موج اجرای کامل، چند نشانهٔ کلیدی قابل ردیابی است. شمارندهٔ برنامه از صفر آغاز می‌شود و در بازه‌های مربوط به مقایسه‌ها به‌صورت پلّکانی بالا می‌رود؛ سپس بسته به بازهٔ ورودی، به یکی از سه آدرس شاخه می‌پرد. این پرش در نما به‌صورت یک جهش ناگهانی در مقدار شمارندهٔ برنامه دیده می‌شود و نه یک پلّهٔ یک‌واحدی.

	اشاره‌گر پشته در سراسر اجرا میان صفر و دو نوسان می‌کند و هرگز از دو بالاتر نمی‌رود؛ همان نتیجه‌ای که در بخش ۱۰.۷ به‌صورت عددی گزارش شد. دو پرچم \lr{S} و \lr{Z} تنها در لحظه‌هایی تغییر می‌کنند که یک \lr{CMP}، \lr{ADD} یا \lr{SUB} اجرا شده باشد و در طول دستورهای \lr{PUSH} و \lr{POP} کاملاً ثابت می‌مانند؛ این همان قید دستورکار است که پرچم‌ها تنها با دستورهای حسابی تغییر کنند.

	سرانجام، خروجی‌های درگاه در انتهای اجرا مقدار می‌گیرند: دو رقم نمایشگر هم‌زمان با نوشتن در \lr{F9} و \lr{FA} به‌روز می‌شوند و چراغ پایان کار در آخرین دستور برنامه روشن می‌گردد. در ورودی‌های خطا، به‌جای آن، چراغ خطا روشن و هر دو رقم نمایشگر صفر می‌مانند.

	\shotw{screenshots/sim_10.png}{0.75}{نمای نزدیک خروجی درگاه‌ها در پایان اجرا}
	\shotw{screenshots/sim_11.png}{0.75}{نمای نزدیک مسیر خطا}

	\section*{۱۰.۱۰ سنتز و تحلیل زمانی}
	\subsection*{الف) مصرف منابع}
	طرح در \lr{Quartus II 13.0.1} و برای خانوادهٔ \lr{Cyclone IV GX} همگردانی شد. انتخاب قطعه روی حالت خودکار بود و ابزار قطعهٔ \lr{EP4CGX15BF14C6} را برگزید. نتیجه \(3{,}208\) عنصر منطقی شامل \(2{,}455\) تابع ترکیبی و \(2{,}162\) ثبات اختصاصی، و ۲۶ پایه است.

	\shotw{screenshots/flow_summary.png}{0.8}{خلاصه‌ی منابع مصرفی سامانه}

	این طرح پرمصرف‌ترین طرح کلّ این مجموعه آزمایش‌هاست و دلیل آن در یک عدد نهفته است: مقدار بیت‌های حافظهٔ مصرف‌شده صفر گزارش شده، در حالی که طرح یک حافظهٔ ۲۵۶ خانه‌ای هشت‌بیتی دارد. یعنی سنتز‌کننده این حافظه را به‌جای یک بلوک حافظهٔ آماده، با \(256\times8=2048\) فلیپ‌فلاپ ساخته است. همین یک آرایه، بیش از نود درصد کلّ \(2{,}162\) ثبات طرح را در بر می‌گیرد و بقیهٔ ثبات‌ها به پشته، شمارندهٔ برنامه، ثبات‌های دستور و عملوند، اشاره‌گر پشته، پرچم‌ها و درگاه‌های خروجی تعلّق دارند.

	دلیل این انتخاب سنتز‌کننده آن است که در این توصیف، حافظه در یک بلوک مقداردهی اوّلیّه با برنامه پر می‌شود و هم‌زمان از دو مسیر ترکیبی خوانده می‌شود؛ چنین الگویی با ساختار بلوک حافظهٔ آمادهٔ قطعه نمی‌خواند. اگر هدف، کاهش مساحت بود، می‌شد حافظه را به‌گونه‌ای بازنویسی کرد که به بلوک حافظهٔ اختصاصی نگاشت شود.

	\subsection*{ب) نماهای ساختاری}
	\shotw{screenshots/rtl_viewer.png}{0.92}{نمای سطح انتقال ثبات‌های سامانه}
	\shotw{screenshots/rtl_top.png}{0.92}{نمای سطح بالای طرح و اتّصال پردازنده به حافظه و درگاه‌ها}
	\shotw{screenshots/post_mapping.png}{0.9}{نمای \lr{Technology Map Viewer (Post-Mapping)}}
	\shotw{screenshots/post_fitting.png}{0.9}{نمای \lr{Technology Map Viewer (Post-Fitting)}}

	\subsection*{ج) تحلیل زمانی}
	برای طرح یک پروندهٔ قیود طراحی\LTRfootnote{Synopsys Design Constraints} نوشته شد که ساعتی با دورهٔ ۲۰ نانوثانیه، معادل ۵۰ مگاهرتز، تعریف می‌کند.

	\docode{code/lab10.sdc}{پرونده‌ی قیود طراحی}

	\shotw{screenshots/timequest.png}{0.85}{خلاصه‌ی تحلیل‌گر زمانی}

	تحلیل‌گر زمانی\LTRfootnote{TimeQuest Timing Analyzer} در بدترین شرایط کاری حاشیهٔ برپایی\LTRfootnote{Slack} \(6.905\) نانوثانیه گزارش می‌کند و مجموع حاشیه‌های منفی صفر است؛ یعنی قید ۵۰ مگاهرتزی با اطمینان برآورده می‌شود. بیشینه‌بسامد\LTRfootnote{Maximum Frequency (Fmax)} برابر \(76.37\,\mathrm{MHz}\) است و بیشینه‌بسامد مقیّد نیز همین مقدار است.

	\shotw{screenshots/fmax_summary.png}{0.75}{گزارش بیشینه‌بسامد سامانه}

	این بیشینه‌بسامد از بیشتر آزمایش‌های پیشین کمتر است و علّت آن همان حافظهٔ ساخته‌شده با ثبات است: خواندن از حافظه یک تسهیم‌کنندهٔ ۲۵۶ به یک می‌سازد که مسیر ترکیبی بلندی دارد و در مسیر بحرانی طرح قرار می‌گیرد.

	با توجّه به برگزاری مجازی آزمایشگاه، ارزیابی این آزمایش بر پایهٔ شبیه‌سازی، سنتز و تحلیل زمانی نرم‌افزاری انجام شده و در پرونده‌های پیکربندی پروژه تخصیص پایه‌ای برای برد ثبت نشده است؛ ازاین‌رو ادّعایی دربارهٔ آزمون فیزیکی روی برد مطرح نمی‌شود.

	\section*{۱۰.۱۱ نتایج}
	\begin{itemize}
		\item پردازندهٔ پشته‌ای با هشت ثبات هشت‌بیتی، هشت دستور پایه و حافظهٔ ۲۵۶ خانه‌ای، مطابق دستورکار پیاده‌سازی شد.
		\item سه دستور \lr{SWAP}، \lr{SHL} و \lr{CMP} افزوده شدند و هر سه با یک میز آزمون جداگانه آزموده گردیدند.
		\item ارتباط با ورودی و خروجی تنها از راه نگاشت در حافظه و با همان دستورهای \lr{PUSH} و \lr{POP} انجام شد.
		\item برنامهٔ زبان ماشین تابع چندبخشی خواسته‌شده را در سه شاخه محاسبه می‌کند و هر دو وضعیّت خطا را تشخیص می‌دهد.
		\item ضرب در دو با \lr{SHL} و ضرب در سه با ترکیب \lr{SHL} و \lr{ADD} ساخته شد، بدون هیچ بلوک ضرب‌کننده‌ای.
		\item رقم پرارزش نمایشگر، در نبود دستورهای منطقی و شیفت راست، با تفریق پیاپی ساخته شد.
		\item بیشینهٔ عمق پشته برای هر ۲۵۶ ورودی ممکن برابر دو اندازه‌گیری شد؛ یعنی یک‌چهارم ظرفیّت پشتهٔ هشت‌تایی.
		\item درستی برنامه پیش از شبیه‌سازی سخت‌افزار، با یک شبیه‌ساز مستقل و برای هر ۲۵۶ ورودی وارسی شد.
		\item سنتز با \(3{,}208\) عنصر منطقی و \(2{,}162\) ثبات انجام شد و تحلیل زمانی حاشیهٔ مثبت \(6.905\) نانوثانیه و بیشینه‌بسامد \(76.37\,\mathrm{MHz}\) داد.
	\end{itemize}

	\section*{نتیجه‌گیری}
	در این آزمایش یک پردازندهٔ کامل، از مجموعه دستور تا برنامهٔ زبان ماشین، ساخته و ارزیابی شد. آنچه این آزمایش را از آزمایش‌های پیشین متمایز می‌کند، دوسویه بودن طراحی است: مجموعه دستور تعیین می‌کند چه برنامه‌هایی نوشتنی‌اند، و نیازهای برنامه تعیین می‌کند چه دستورهایی باید افزوده شوند. هر سه دستور تازه از همین جهت زاده شدند؛ \lr{SHL} چون ضرب لازم بود، و \lr{CMP} چون مقایسه بدون از دست دادن عملوند لازم بود. محدودیّت‌های مجموعه دستور نیز خود را در جاهای دیگری نشان دادند: نبود دستور شیفت راست، ساختن رقم پرارزش نمایشگر را به یک حلقهٔ تفریق پیاپی تبدیل کرد. سرانجام، گزارش سنتز یادآور شد که سادگی یک معماری به معنای ارزان بودن پیاده‌سازی آن نیست؛ حافظهٔ ۲۵۶ خانه‌ای که در سطح توصیف تنها یک آرایه است، در سخت‌افزار بیش از دو هزار فلیپ‌فلاپ گرفت و مسیر بحرانی طرح را نیز همان ساخت.
\end{document}
